\reading{IM-14}{Madoff: A Riot of Red Flags}{STRIKE FIRST}{3}

\srcnote{Greg Gregoriou and Francois Lhabitant, \textit{Madoff: A Riot of Red
Flags} (EDHEC Publication, 2009). No GARP source book covers this reading; the
material comes from the compiled notes supplied with the source set, one rung
below a GARP source.}

\begin{imkeybox}[title={The two objectives in this reading}]
\begin{enumerate}[label=\textbf{\alph*.}]
\item Explain the history of Bernie Madoff's hedge fund and describe the fund's reported strategy and its reported performance.
\item Identify characteristics of Bernie Madoff's hedge fund that can serve as warning signs or potential ``red flags'' and explain how they can indicate potential fraud risk.
\end{enumerate}
\end{imkeybox}

% ------------------------------------------------------------------
\lo{IM-14 a}{Explain the history of Bernie Madoff's hedge fund and describe the fund's reported strategy and its reported performance.}

Bernard L. Madoff Investment Securities (BLMIS) was founded in \term{1960}. The
market-making and broker-dealer arm was a \term{genuine and successful business},
an early adopter of electronic trading and at times among the largest market
makers in NYSE-listed stock traded off-exchange. Madoff held senior industry
positions, including \term{non-executive chairman of NASDAQ} in the early 1990s,
and served on advisory committees.

\begin{imtrapbox}[breakable=false,title={The respectability was part of the cover}]
Few were willing to believe an establishment figure could be running a Ponzi
scheme --- and the legitimate broker-dealer business gave the fraudulent advisory
business its credibility. A question that describes BLMIS as wholly fraudulent has
missed the structure: \term{one arm was real}, and that is what made the other
one survivable for decades.
\end{imtrapbox}

\subsection{The reported performance}

\begin{itemize}
\item Returns \term{far too smooth}, with almost no down months.
\item \term{Near-zero correlation to equity markets}, despite a portfolio
explicitly built from S\&P 100 stocks.
\item An implied \term{Sharpe ratio far beyond anything the strategy could
generate}. Equity-like returns with bond-like volatility from a long-equity
collar is not a difficult feat; it is an impossible one.
\end{itemize}

\begin{imkeybox}
\term{The shape of the return series was the single most diagnostic piece of
evidence}, and it was available to anyone from the published numbers alone.
Nothing confidential was required to see it.
\end{imkeybox}

% ------------------------------------------------------------------
\lo{IM-14 b}{Identify characteristics of Bernie Madoff's hedge fund that can serve as warning signs or potential ``red flags'' and explain how they can indicate potential fraud risk.}

\renewcommand{\arraystretch}{1.25}
\noindent\begin{tabularx}{\textwidth}{@{}>{\raggedright\arraybackslash}p{2.2cm}XX@{}}
\toprule
\textbf{\sffamily\small Category} & \textbf{\sffamily\small Red flag} & \textbf{\sffamily\small Why it signals fraud risk} \\
\midrule
Returns and strategy & Returns far too smooth; almost no down months. Near-zero correlation to equities despite \emph{holding} equities. Performance unaffected by regime. & A return series that does not respond to the market it is invested in is not a strategy outcome. It is a number being chosen. \\
Audit & \term{Friehling \& Horowitz}, a three-person firm operating from a small suburban office, audited a multi-billion-dollar operation. The firm had represented to the AICPA \emph{for years} that it did not conduct audits. & \term{Audit capacity must be proportionate to the entity.} A manifestly undersized auditor cannot provide assurance, and suggests one has been chosen \emph{for} that reason. \\
Governance & Family members in key compliance and operations roles. Extreme secrecy about method. Investors who asked detailed questions were \term{threatened with having their accounts returned}. & Independent compliance cannot exist where the compliance officer reports to a relative. \term{Punishing due diligence inverts the normal relationship} and is a serious warning in itself. \\
Economics of the fee & \term{No management fee and no performance fee} --- compensation taken instead through trading commissions on the broker-dealer side. & A manager forgoing the industry's standard economics is either uneconomic or is being paid another way. Either answer demands an explanation. \\
\bottomrule
\end{tabularx}
\renewcommand{\arraystretch}{1}
\figcap{Four categories, and every one of them visible from outside. This is why
the reading sits in the curriculum next to IM-12: the fee structure, the auditor
and the governance are all items on a due diligence questionnaire. Note the
inversion in the governance row --- a manager who \emph{penalises} the investor
for asking is the clearest of the four signals and the hardest to rationalise.}

\begin{imtrapbox}[title={Consolidated trap summary --- IM-14}]
\begin{itemize}
\item \term{Scope, IM-14 a.} The broker-dealer arm was genuine. The fraud was in
the advisory business.
\item \term{Definition, IM-14 a.} The diagnostic evidence was the \emph{shape of
the return series} --- smoothness, absent drawdowns, near-zero correlation to the
equities it held --- not any confidential document.
\item \term{Role, IM-14 b.} The auditor red flag is about \emph{capacity
proportionate to the entity}, not about the auditor's reputation.
\item \term{Sign, IM-14 b.} \term{No fees} was a red flag, not a benefit.
Compensation came through commissions instead.
\item \term{Polarity, IM-14 b.} Investors who asked questions were threatened with
redemption. A manager punishing due diligence is itself the signal.
\end{itemize}
\end{imtrapbox}

% ==================================================================
